\documentclass[10pt,a4paper]{amsart}
\usepackage[margin=19mm]{geometry}
\usepackage{amsmath,amssymb,mathtools}
\usepackage[scr=rsfs,cal=euler]{mathalfa}
\usepackage{xfrac}
\usepackage[unicode,hidelinks,bookmarks=false]{hyperref}
\newcommand{\ie}{i.e.\ }
\newcommand{\cf}{cf.\ }
\newcommand{\resp}{respectively\ }
\newcommand{\Subsection}{\subsection}
\providecommand{\nobreakdash}{\nobreak\hbox{-}}
\setlength{\emergencystretch}{2em}
\multlinegap0pt
\usepackage{bm}
\usepackage{bbm}
\usepackage{dsfont}
\usepackage{xspace}
\usepackage{cancel}
\usepackage{mathtools}
\usepackage{amsthm}
\makeatletter
\@ifundefined{theorem}{\newtheorem{theorem}{Theorem}}{}
\@ifundefined{lemma}{\newtheorem{lemma}[theorem]{Lemma}}{}
\makeatother
\newcommand{\rk}{\textnormal{rk}}
\newcommand{\srk}{\textnormal{srk}}
\renewcommand{\P}{\mathop{\mathbb{P}}}
\renewcommand{\k}{{\mathbf{k}}}
\title[Slice rank and analytic rank for trilinear forms]{Slice rank and analytic rank for trilinear forms}
\author{Amichai Lampert}
\date{Source: arXiv:2404.19704v2 (CC BY 4.0)}
\begin{document}
\maketitle

\noindent\emph{Source and licence.} Slice rank and analytic rank for trilinear forms. Version arXiv:2404.19704v2 (CC BY 4.0), distributed under the Creative Commons Attribution 4.0 International licence, as recorded in the arXiv licence metadata for this exact version and shown on the abstract page https://arxiv.org/abs/2404.19704v2. Full rights notice: licenses/arxiv-2404.19704-CC-BY-4.0.txt.\par\smallskip

\noindent\textit{Editorial scope.} Counted unit: the Lemma whose source label is \texttt{low-rk-der} in the article (source0.tex lines 246--249, source label \texttt{low-rk-der}), delivered together with the article's own complete original proof of it (source0.tex lines 251--289, one \texttt{proof} environment of 39 lines). No mathematics is written, repaired or re-derived: statement and proof are the article's own text line for line. Required context retained uncounted: none. Every definition, display and label of those lines is retained unchanged. Omitted from this package and not redistributed: 1--245, 250--250, 290--414. Nothing inside a retained excerpt is omitted or altered: every retained body is delivered exactly as the article's lines are written, so no omission receipt of any kind accompanies this package. Citations: the retained excerpts carry no citation command at all, so no bibliography entry of the article is transcribed and no reference is redistributed. Reference adaptation: none was needed and none was made; every reference inside the delivered unit resolves inside it, so no unresolved reference or citation is produced. Printed slips kept literal: no printed word, symbol, hypothesis or number of the counted statement or of its proof was altered. Layout and class: the article's own class is \texttt{daj} with packages including amscd, amssymb, amsmath, amsfonts, amsthm, comment, hyperref, xcolor; the wrapper is the lane's pinned \texttt{amsart} editorial wrapper at 10pt on A4, which restates the theorem-like environments the retained text needs and adds the article's own macro definitions that the retained text needs (\texttt{\textbackslash P}, \texttt{\textbackslash k}, \texttt{\textbackslash rk}, \texttt{\textbackslash srk}), with extra wrapper packages (bm, bbm, dsfont, xspace, cancel, mathtools). The delivered numbering is the unit's own. Assets: no retained span contains a figure environment, a graphics-inclusion command, a PDF-page inclusion, a table or a TikZ picture; no image file is copied into this package, the package declares no asset dependency and no image file is redistributed with it. Licence: version arXiv:2404.19704v2 of this article is distributed under the Creative Commons Attribution 4.0 International licence, as recorded in the arXiv licence metadata for this exact version and shown on the abstract page https://arxiv.org/abs/2404.19704v2. A full rights notice is delivered at licenses/arxiv-2404.19704-CC-BY-4.0.txt. Characters: every retained excerpt is pure ASCII, so the delivered engine, XeLaTeX with its default Latin Modern text font, reports no missing character.
\begin{lemma}\label{low-rk-der}
    If $g:U\times V\times W\to k$ is a trilinear form with
    $$\P_{u\in U} (\rk g[u] > t) < \frac{q-1}{2qt}$$ for some positive integer $t$ then $\srk (g) \le 4t.$ Moreover, the linear forms in the slice rank decomposition are obtained by fixing coordinates of $g.$
    \end{lemma}

    \begin{proof}
         Let $A = \{u\in U: \rk g[u] \le t \}.$ Our assumption is that $\frac{|A|}{|U|} > 1-\frac{q-1}{2qt}.$ The proof proceeds by induction on $t.$
         
         \textbf{Base case $t=1:$} Suppose first that $g[u] = 0$ for all $u\in A.$ Since $|A| > |U|/q$ and $u\mapsto g[u]$ is a linear map, we deduce that $g[u] = 0$ for all $u\in U,$ so $g = 0.$  We may henceforth assume that there is some $u_0\in A$ with $\rk g[u_0] = 1$. Writing $g[u_0] = \alpha(v)\beta(w),$ we claim that for all $u\in A\cap (A-u_0)$ the bilinear form $g[u]$ is contained in the ideal $(\alpha,\beta).$  Indeed, if $g[u] = \gamma(v)\delta(w)$ with $\gamma \not\in \text{span}(\alpha)$ and $\delta\not\in\text{span}(\beta)$ then 
         \[
         \rk g[u_0+u] = \rk(\alpha(v)\beta(w)+\gamma(v)\delta(w)) = 2,
         \]
         contradicting the fact that $u_0+u \in A.$ Therefore, if $V_0 = \{v: \alpha(v) = 0\}$ and $W_0 = \{w: \beta(w) = 0\}$ and $\tilde g = g\restriction_{U\times V_0\times W_0},$ we get that $\tilde g[u] = 0$ whenever $u\in A\cap (A-u_0),$  a set of density greater than $ 1- 2\cdot \frac{q-1}{2q} = \frac{1}{q}.$ By the linearity of $u\mapsto \tilde g[u]$ we deduce $\tilde g = 0$, so $g\in(\alpha,\beta)\implies \srk(g) \le 2.$ 

         
         \textbf{Inductive step:} We may assume that there is some $u_0\in A$ with $\rk g[u_0] = t,$ otherwise we're done by the inductive hypothesis. Write $g[u_0] = \sum_{i=1}^t \alpha_i(v)\beta_i(w)$ and set $V_0 = \{v:\alpha_i(v) = 0\ \forall i\in[t]\},$ $W_0 = \{w:\beta_i(w) = 0\ \forall i\in[t]\}.$ We claim that $\tilde g = g\restriction_{U\times V_0\times W_0}$ satisfies $\rk \tilde g[u] \le t/2 $ for all $u\in A\cap (A-u_0).$ This will complete the proof because this set of $u$'s has density greater than $ 1-2\cdot \frac{q-1}{2qt} \ge 1- \frac{q-1}{2q\lfloor t/2\rfloor}, $ so the inductive hypothesis implies 
         \[
         \srk (\tilde g) \le 4\lfloor t/2\rfloor \le 2t \implies \srk(g) \le 4t.
         \]
        To prove the claim about $\rk \tilde g[u],$ suppose it has rank $s$ and write $\tilde g[u] = \sum_{i=1}^s \tilde\gamma_i(v) \tilde \delta_i(w).$ Let $\gamma_i,\delta_i$ be linear forms on $V,W$ which restrict to $\tilde\gamma_i,\tilde\delta_i$, respectively.  Note that our assumptions on $\rk g[u_0]$ and $\rk \tilde g[u]$ imply that  both $\alpha_i,\gamma_i$ and  $\beta_i,\delta_i$ are linearly independent. Choosing bases, we may identify $V=\k^n,W=\k^m$ and 
        $$ \gamma_i(v) = v_i,\ \delta_i(w) = w_i,\  \alpha_i(v) = v_{s+i},\ \beta_i(w) = w_{s+i}.$$ This yields an equation
         \[
         g[u] = \sum_{i=1}^s v_iw_i +\sum_{j=1}^t  v_{s+j} \phi_j (w)+\sum_{j=1}^t  w_{s+j} \tau_j (v).
         \]
         Decomposing $$\phi_j(w) = \sum_{i=1}^s a_{i,j} w_i + \phi'(w_{s+1},\ldots,w_m),\ \tau_j(v) = \sum_{i=1}^s b_{i,j} v_i + \tau'(v_{s+1},\ldots,v_n),$$ 
         we get 
         \[
         g[u] = \sum_{i=1}^s  \gamma'_i(v) \delta'_i(w) +q(v_{s+1},\ldots,v_n,w_{s+1},\ldots,w_m),
         \]
         where 
         \[
         \gamma'_i(v) = v_i+\sum_{j=1}^t a_{i,j} v_{s+j},\ \delta'_i(w)= w_i+\sum_{j=1}^t b_{i,j} w_{s+j}.
         \]
         The collection of linear forms $\gamma'_1,\ldots,\gamma'_s,v_{s+1},\ldots,v_n$ spans $V^*$ and so is linearly independent, likewise for $\delta'_1,\ldots,\delta'_s,w_{s+1},\ldots,w_m.$ This means that 
         \begin{equation}\label{rk-q}
             t\ge \rk g[u] = s+ \rk(q).
         \end{equation}
         Since $u+u_0\in A,$ we get 
         \begin{align*}
             t &\ge \rk(g[u]+g[u_0]) = \rk \left(\sum_{i=1}^s  \gamma'_i(v) \delta'_i(w)+\sum_{i=1}^t v_{s+i} w_{s+i}+q\right) \\
             &= s+\rk\left(\sum_{i=1}^t v_{s+i} w_{s+i}+q\right) \ge s+t-\rk(q) \ge 2s
         \end{align*}
         where the last inequality used \eqref{rk-q}. This proves that $\rk\tilde g[u] = s\le t/2,$ completing the proof of the lemma.
    \end{proof}

\end{document}
